% !TEX program = platex
% 数学ノート集・共通テンプレート v2（積分論2026年度第1部の版面を参考）
% A5判。章内で定義・定理・例・注意等の番号を共通化。
% コンパイル： platex -interaction=nonstopmode math_note_template_v2.tex (2回)
%              dvipdfmx math_note_template_v2.dvi
\documentclass[a5paper,10pt,onecolumn]{jsarticle}
\usepackage[dvipdfmx,top=22mm,bottom=20mm,inner=19mm,outer=18mm,headheight=13pt,headsep=6mm]{geometry}
\usepackage{amsmath,amssymb,amsthm,mathtools}
\usepackage{enumitem,fancyhdr,titlesec}
\usepackage[dvipdfmx,bookmarks=true,colorlinks=false]{hyperref}
\usepackage{pxjahyper}

% ----- 差し替える文書情報 -----
\newcommand{\NoteTitle}{Weyl代数と微分作用素環}
\newcommand{\NoteSeries}{数学ノート集}

% ----- 文字・組版 -----
% 参考PDFの質素な明朝＋Computer Modernの印象を引き継ぎながら，
% A5判で窮屈にならないよう本文の行間を少し確保する．
\linespread{1.13}
\setlength{\parindent}{1zw}
\setlength{\parskip}{0pt}
\setlength{\abovedisplayskip}{6pt plus 2pt minus 2pt}
\setlength{\belowdisplayskip}{6pt plus 2pt minus 2pt}
\setlength{\abovedisplayshortskip}{3pt plus 2pt minus 1pt}
\setlength{\belowdisplayshortskip}{4pt plus 2pt minus 1pt}
\setlist[enumerate]{topsep=3pt,itemsep=1pt,parsep=0pt,leftmargin=2.4zw}
\titleformat{\section}{\normalfont\normalsize\bfseries}{\thesection}{0.8zw}{}
\titlespacing*{\section}{0pt}{2.0ex plus .4ex}{.9ex}
\titleformat{\subsection}{\normalfont\normalsize\bfseries}{\thesubsection}{0.7zw}{}
\titlespacing*{\subsection}{0pt}{1.5ex plus .3ex}{.6ex}
\raggedbottom

% ----- 柱とページ番号：ページ上部、罫線なし -----
\fancypagestyle{mathnote}{%
  \fancyhf{}%
  \fancyhead[L]{\footnotesize \NoteSeries\ （\NoteTitle）}%
  \fancyhead[R]{\footnotesize \thepage}%
  \renewcommand{\headrulewidth}{0pt}%
  \renewcommand{\footrulewidth}{0pt}%
}
\pagestyle{mathnote}

% ----- 節内通し番号の定義・定理群 -----
\newtheoremstyle{mathnotetheorem}
 {4.5pt}{4.5pt}{\normalfont}{}%
 {\bfseries}{}{0.65zw}%
 {\thmname{#1}\thmnumber{ #2}\thmnote{（#3）}}
\theoremstyle{mathnotetheorem}
\newtheorem{theorem}{定理}[section]
\newtheorem{definition}[theorem]{定義}
\newtheorem{lemma}[theorem]{補題}
\newtheorem{proposition}[theorem]{命題}
\newtheorem{corollary}[theorem]{系}
\newtheorem{example}[theorem]{例}
\newtheorem{remark}[theorem]{注意}
\renewcommand{\proofname}{証明}
\renewcommand{\qedsymbol}{$\square$}
\makeatletter
\renewenvironment{proof}[1][\proofname]{%
  \par\pushQED{\qed}\normalfont\topsep4\p@\@plus2\p@\relax
  \trivlist\item[\hskip\labelsep\bfseries #1]\ignorespaces
}{\popQED\endtrivlist\@endpefalse}
\makeatother

% ----- 使うノートで適宜拡張する記号 -----
\newcommand{\R}{\mathbb{R}}
\newcommand{\norm}[1]{\lVert #1\rVert}
\newcommand{\abs}[1]{\lvert #1\rvert}

\begin{document}
\section*{読了後にできること}
本稿を読むための前提は，非可換環，微分作用素，交換子．
本稿の具体的な到達課題は，Weyl代数で微分演算子を正規形へ並べ替える．

\section{Weyl代数の定義}
本稿では $k$ を標数零の体とする．
\begin{definition}
$n$ 次Weyl代数 $A_n(k)$ とは，単位元を持つ結合的
$k$-代数で，生成元
$x_1,\ldots,x_n,\partial_1,\ldots,\partial_n$ が
\[
 x_ix_j=x_jx_i,\qquad
 \partial_i\partial_j=\partial_j\partial_i,\qquad
 \partial_ix_j-x_j\partial_i=\delta_{ij}
 \tag{1}
\]
を満たす自由結合代数の商として定義されるものをいう．
\end{definition}
\begin{lemma}[微分作用素による表現]\label{lem:operators}
$k[z_1,\ldots,z_n]$ 上で $x_i$ を
$z_i$ の乗法，$\partial_i$ を $z_i$ に関する
偏微分として作用させると，式 (1) の関係が成立する．
従ってWeyl代数から
$\operatorname{End}_k(k[z_1,\ldots,z_n])$ への
$k$-代数準同型が定まる．
\end{lemma}
\begin{proof}
乗法作用素は互いに可換であり，
偏微分作用素も互いに可換である．
任意の多項式 $f$ に対しLeibniz則から
\[
 \partial_i(z_jf)-z_j\partial_if=\delta_{ij}f
\]
となる．従って定義関係が全て成り立ち，
商代数の普遍性によって準同型が誘導される．
\end{proof}
\begin{remark}
標数零の仮定は，後の標準形の一次独立性および中心の記述，
有限次元表現の不存在において本質的である．
\end{remark}

\section{正規形とPBW型基底}
多重指数 $\alpha,\beta\in\mathbb N^n$ に対し
$x^\alpha=x_1^{\alpha_1}\cdots x_n^{\alpha_n}$，
$\partial^\beta=\partial_1^{\beta_1}\cdots\partial_n^{\beta_n}$
と書く．
\begin{lemma}[正規順序への書換え]\label{lem:normal}
$A_n(k)$ の任意の元は
$x^\alpha\partial^\beta$ の有限個の $k$-線形結合で表される．
\end{lemma}
\begin{proof}
生成元からなる語を考える．
全ての $x_i$ を左，$\partial_j$ を右に移すように
隣接する $\partial_i x_j$ を交換する．
関係 (1) から
$\partial_i x_j=x_j\partial_i+\delta_{ij}$ である．
第1項では「左の微分記号と右の乗法記号」の逆転数が
一つ減り，第2項は語の長さが二つ減る．
語の長さと逆転数に関する辞書式帰納法を用いれば，
この書換えは有限回で終了する．
$x_i$ 同士，$\partial_j$ 同士は可換なので，
得られた語は $x^\alpha\partial^\beta$ の形に整理できる．
\end{proof}
\begin{theorem}[Weyl代数のPBW型基底]\label{thm:pbw}
\[
 \{x^\alpha\partial^\beta:
   \alpha,\beta\in\mathbb N^n\}
\]
は $A_n(k)$ の $k$-線形基底である．
特に補題\ref{lem:operators}の表現は忠実である．
\end{theorem}

\begin{proof}
張ることは補題\ref{lem:normal}による．
一次独立性を示すため，有限和
$P=\sum_{\alpha,\beta}
 c_{\alpha\beta}x^\alpha\partial^\beta$ が
多項式環上の作用素として零と仮定する．
形式的冪級数環
$k[z_1,\ldots,z_n][[t_1,\ldots,t_n]]$ で
\[
 E(z,t)=\exp(t_1z_1+\cdots+t_nz_n)
 =\sum_{\nu\in\mathbb N^n}{t^\nu z^\nu\over\nu!}
\]
を考える．標数零なので各 $\nu!$ は可逆である．
$P$ は $z$ のみに作用する有限階微分作用素だから，
$t$ に関する各係数へ作用させることで
形式的冪級数全体に拡張できる．
$P$ は全ての多項式上で零なので $PE=0$ である．
一方 $\partial_iE=t_iE$ より
\[
 PE=\left(\sum_{\alpha,\beta}
 c_{\alpha\beta}z^\alpha t^\beta\right)E.
\]
$E$ の定数項は1なので冪級数として可逆である．
従って $\sum_{\alpha,\beta}c_{\alpha\beta}
z^\alpha t^\beta=0$ が成立する．
多変数多項式の単項式の一次独立性から
全ての $c_{\alpha\beta}=0$ である．
ゆえに元の正規形は一次独立であり，
同時に作用素表現の核も零である．
\end{proof}
\begin{proposition}[一変数の交換公式]
$A_1(k)$ で $a,b\ge0$ とすると
\[
 \partial x^a=x^a\partial+a x^{a-1}
 \quad(a\ge1),
\]
および
\[
 \partial^b x^a
 =\sum_{j=0}^{\min(a,b)}
 {b\choose j}a(a-1)\cdots(a-j+1)
 x^{a-j}\partial^{b-j}
 \tag{2}
\]
が成立する．$j=0$ の積は1とする．
\end{proposition}
\begin{proof}
第一式は $\partial x=x\partial+1$ から
$a$ に関する帰納法で得られる．
第二式は $b=0$ では自明である．
$b$ について成立すると仮定して左から $\partial$ を掛け，
第一式で正規順序へ直す．
$x^{a-j}\partial^{b-j+1}$ の項と
$(a-j)x^{a-j-1}\partial^{b-j}$ の項が現れる．
同じ単項式の係数をまとめると
$ {b\choose j}+{b\choose j-1}={b+1\choose j}$
によって $b+1$ の場合の式を得る．
\end{proof}

\section{次数フィルトレーション}
\begin{definition}
$m\ge0$ に対して
\[
 F_mA_n=\operatorname{span}_k
 \{x^\alpha\partial^\beta:|\alpha|+|\beta|\le m\}
\]
と置き，$F_{-1}A_n=0$ とする．
\[
 \operatorname{gr}A_n=
 \bigoplus_{m\ge0}F_mA_n/F_{m-1}A_n
\]
を付随次数環という．
\end{definition}
\begin{lemma}
$F_pA_n\cdot F_qA_n\subset F_{p+q}A_n$ である．
さらに $u\in F_pA_n$，$v\in F_qA_n$ なら
$[u,v]\in F_{p+q-2}A_n$ である
（負の添字の $F_m$ は零とする）．
\end{lemma}
\begin{proof}
正規形の二単項式を掛け，関係 (1) に従って
乗法記号を左へ移すと，
交換しない項では全次数は $p+q$ 以下に保たれる．
交換によって現れる定数項は記号二個を消すので，
その項は次数を二つ下げる．
交換子を取ると最高次数の項は可換積として一致して相殺する．
従って交換子には少なくとも一度の交換に由来する項しか残らず，
その次数は高々 $p+q-2$ である．
一般の $u,v$ は有限な線形結合なので結論が従う．
\end{proof}
\begin{theorem}[付随次数環]\label{thm:gr}
\[
 \operatorname{gr}A_n(k)\cong
 k[X_1,\ldots,X_n,\Xi_1,\ldots,\Xi_n]
\]
が次数付き $k$-代数として成立する．
\end{theorem}
\begin{proof}
前補題から付随次数環は可換であり，
$x_i,\partial_j$ の主記号をそれぞれ
$X_i,\Xi_j$ に対応させることで
右辺から左辺への次数付き準同型が得られる．
PBW型基底定理\ref{thm:pbw}によって
次数 $m$ の成分では
$|\alpha|+|\beta|=m$ の主記号が $k$-線形基底となる．
右辺の同次数単項式も同じ添字を持つ基底であるから，
準同型は各次数で同型であり，全体でも同型である．
\end{proof}

\section{Weyl代数のNoether性}
\begin{lemma}[付随次数環からの有限生成性]\label{lem:filtered}
$A$ を非負整数で添字付けられた増大フィルトレーション
$F_0A\subset F_1A\subset\cdots$ を持つ環とし，
$A=\bigcup_mF_mA$，
$F_pA\,F_qA\subset F_{p+q}A$ を仮定する．
$\operatorname{gr}A$ が左Noether環なら $A$ は左Noether環である．
\end{lemma}
\begin{proof}
$A$ の任意の左イデアル $I$ に対し
\[
 \operatorname{gr}I=
 \bigoplus_{m\ge0}(I\cap F_mA)/(I\cap F_{m-1}A)
\]
を考える．これは $\operatorname{gr}A$ の次数付き左イデアルである．
Noether性から斉次元である有限個の生成元
$\sigma(g_1),\ldots,\sigma(g_s)$ を取れる．
各生成元は $g_i\in I$ の主記号として選んでよい．
任意の零でない $f\in I$ のフィルトレーション次数を $m$ とする．
$\sigma(f)$ はこれらの生成元の斉次な積の和
$\sum_i a_i^\#\sigma(g_i)$ と書けるので，
各 $a_i^\#$ を適当な $a_i\in A$ に持ち上げる．
このとき
$f-\sum_i a_i g_i$ は $I$ に属し，
その次数は $m$ より真に小さい．
非負整数の降下帰納法によって
$f-\sum_i a_i g_i$ は $(g_1,\ldots,g_s)$ に属する．
従って $f$ もこの左イデアルに属する．
任意の $f\in I$ に対して成立するので
$I=Ag_1+\cdots+Ag_s$ は有限生成である．
\end{proof}
\begin{theorem}[Weyl代数の両側Noether性]\label{thm:noether}
$A_n(k)$ は左Noether環かつ右Noether環である．
\end{theorem}
\begin{proof}
定理\ref{thm:gr}から
$\operatorname{gr}A_n$ は有限変数多項式環である．
これはHilbertの基底定理によってNoether環である．
補題\ref{lem:filtered}から $A_n$ は左Noetherとなる．
右Noether性には反自己同型
$\tau:A_n\to A_n$ を使う．
生成元で
$\tau(x_i)=x_i$，
$\tau(\partial_i)=-\partial_i$ と置き，
$\tau(ab)=\tau(b)\tau(a)$ として延長する．
例えば
\[
 \tau(\partial_ix_j-x_j\partial_i)
 =-x_j\partial_i+\partial_ix_j=\delta_{ij}
\]
なので関係式を保つ．
また $\tau^2=\mathrm{id}$ である．
従って $\tau$ は左イデアルと右イデアルを互いに対応させ，
有限生成性を保つ．
左Noether性から右Noether性が従う．
\end{proof}
\section{中心と有限次元表現}
\begin{theorem}[Weyl代数の中心]\label{thm:center}
標数零の体 $k$ 上で
$Z(A_n(k))=k$ である．
\end{theorem}
\begin{proof}
$P=\sum_{\alpha,\beta}
c_{\alpha\beta}x^\alpha\partial^\beta$ が中心に属するとする．
生成元 $x_j$ との交換子を計算すると
\[
 [P,x_j]=
 \sum_{\alpha,\beta}
 c_{\alpha\beta}\beta_j
 x^\alpha\partial^{\beta-e_j}=0
\]
となる（$\beta_j=0$ の項は零とする）．
PBW基底の一次独立性と標数零から
$\beta_j>0$ の係数は全て零である．
全ての $j$ にこれを適用すると $P\in k[x_1,\ldots,x_n]$．
さらに
$[P,\partial_j]=-\partial P/\partial x_j=0$ なので，
再び標数零を用いれば $P$ は定数である．
逆に体 $k$ の元は全ての元と可換なので中心に含まれる．
\end{proof}
\begin{theorem}[非零有限次元加群の不存在]
標数零の体 $k$ 上の $A_n(k)$ は，
$n\ge1$ のとき非零の有限次元左加群を持たない．
\end{theorem}
\begin{proof}
有限次元左加群 $V$ が存在すると仮定し，
$\partial_1,x_1$ の作用をそれぞれ線形写像 $D,X$ とする．
定義関係から
$DX-XD=\mathrm{id}_V$ が成立する．
有限次元空間で
$\operatorname{tr}(DX)=\operatorname{tr}(XD)$ だから
左辺の跡は零である．
一方右辺の跡は $\dim_kV$ であり，
$V\ne0$ なら正整数だから標数零の体で非零である．
矛盾するので $V=0$ しか存在しない．
\end{proof}

\section{微分方程式からD加群へ}
\begin{definition}
$A_1(k)$ の元 $P$ に対して左巡回加群
\[
 M_P=A_1(k)/A_1(k)P
\]
を考える．その標準生成元を $e=1+A_1(k)P$ とすると
$Pe=0$ が成立する．
\end{definition}
\begin{example}[一次線形微分方程式]
$\lambda\in k$ とし
$P=\partial-\lambda$ と置く．
$M_P$ は $k[z]$ と $k$-ベクトル空間として同型であり，
その上で $x$ は $z$ の乗法，
$\partial$ は $d/dz+\lambda$ として作用する．
\end{example}
\begin{proof}
$k[z]$ 上に $x\cdot f=zf$，
$\partial\cdot f=f'+\lambda f$ と定めると
\[
 (\partial x-x\partial)f
 =(zf)'+\lambda zf-z(f'+\lambda f)=f
\]
なのでWeyl代数の左加群となる．
この加群は $1$ を $x$ の作用で繰り返して生成でき，
$(\partial-\lambda)\cdot1=0$ である．
従って $M_P$ から $k[z]$ への全射準同型が得られる．
一方PBW基底で表した各元の $\partial$ の冪を
$\partial=(\partial-\lambda)+\lambda$ に置き換え，
右端の $(\partial-\lambda)$ を法として消去すると，
商 $M_P$ の各元は $x^a e$ の有限線形結合で表される．
これらの像は $z^a$ で一次独立なので，準同型は単射でもある．
\end{proof}
\begin{proposition}[解空間の具体例]
$k=\mathbb C$ とし，複素平面上の正則関数 $u$ について
$(d/dz-\lambda)u=0$ が成立することと，
ある定数 $C\in\mathbb C$ に対し
$u(z)=Ce^{\lambda z}$ となることは同値である．
\end{proposition}
\begin{proof}
$(e^{-\lambda z}u(z))'
 =e^{-\lambda z}(u'(z)-\lambda u(z))$ なので，
微分方程式が成り立つと
$e^{-\lambda z}u(z)$ は全平面で導関数が零となる．
連結な領域上の正則関数で導関数が零なら定数だから
$u(z)=Ce^{\lambda z}$ である．
逆向きは直接微分すればよい．
\end{proof}
\begin{remark}
多変数の微分方程式系を微分作用素環の左イデアルで表し，
その商左加群として扱うことがD加群論の出発点である．
良いフィルトレーション，特性多様体，正則ホロノミーなどは
後続の独立したノートで扱う．
\end{remark}

\section{到達課題の完全解答}
\noindent\textbf{Weyl代数の交換関係を計算する．}
\begin{proposition}
標数零の体 $k$ 上のWeyl代数
$A_1(k)=k\langle x,\partial\rangle/
(\partial x-x\partial-1)$ において，
$n\ge1$ なら
\[
\partial x^n=x^n\partial+nx^{n-1}
\]
である．
\end{proposition}
\begin{proof}
$n=1$ は定義関係である．
$n$ について成立すると仮定すると，
\[
\begin{aligned}
\partial x^{n+1}
&=(\partial x^n)x\\
&=x^n\partial x+nx^n\\
&=x^{n+1}\partial+(n+1)x^n
\end{aligned}
\]
となる．帰納法により成立する．
\end{proof}
\begin{corollary}
$[\partial,x^n]=nx^{n-1}$ であり，
多項式 $f(x)$ に対し
$[\partial,f(x)]=f'(x)$ が成立する．
\end{corollary}
\begin{proof}
第一式は上式から $x^n\partial$ を引けばよい．
第二式は多項式の線形性による．
\end{proof}
\noindent\textbf{到達確認．}
$\partial^2x^2$ を $x$ を左にした正規形にせよ．
\par\noindent\textbf{解答．}
$\partial x^2=x^2\partial+2x$ を使って
$\partial^2x^2
=\partial(x^2\partial+2x)
=x^2\partial^2+4x\partial+2$．

\section{正規形算法と主記号の具体計算}
\begin{proposition}[一変数での積の正規形]
$P=x^a\partial^b$，$Q=x^c\partial^d$ とする．
このとき
\[
 PQ=\sum_{j=0}^{\min(b,c)}
 \binom bj c^{\underline j}
 x^{a+c-j}\partial^{b+d-j},
\]
ただし $c^{\underline 0}=1$，
$c^{\underline j}=c(c-1)\cdots(c-j+1)$ とする．
\end{proposition}
\begin{proof}
既に示した交換公式 (2) を
$x^a(\partial^b x^c)\partial^d$ に代入する．
$x$ 同士および $\partial$ 同士を整理すれば
表示式を得る．
\end{proof}
\begin{example}[演算子の積と交換子]
$P=x^2\partial$，$Q=x\partial^2$ に対して
\[
 \begin{aligned}
 PQ&=x^2(\partial x)\partial^2
     =x^3\partial^3+x^2\partial^2,\\
 QP&=x(\partial^2x^2)\partial
     =x^3\partial^3+4x^2\partial^2+2x\partial.
 \end{aligned}
\]
従って
\[
 [P,Q]=-3x^2\partial^2-2x\partial.
\]
この交換子の全次数は4であり，
$P,Q$ の全次数の和6より2低い．
\end{example}
\begin{proposition}[主記号と交換子]
$P,Q\in A_1(k)$ の全次数をそれぞれ $p,q$ とする．
$[P,Q]$ の次数 $p+q-2$ の斉次部分の主記号は
\[
 \frac{\partial\sigma_p(P)}{\partial\Xi}
 \frac{\partial\sigma_q(Q)}{\partial X}
 -
 \frac{\partial\sigma_p(P)}{\partial X}
 \frac{\partial\sigma_q(Q)}{\partial\Xi}
\]
である．ここで $[P,Q]$ の次数が $p+q-2$
未満の場合，左辺は零と解釈する．
\end{proposition}
\begin{proof}
双線形性により $P=x^a\partial^b$，
$Q=x^c\partial^d$ の場合を示せば十分である．
積の公式で $j=0$ の項は $PQ$ と $QP$ で一致する．
$j=1$ の係数は $bc$ と $da$ であり，
次数は $a+b+c+d-2$ である．
$j\ge2$ の項は少なくとも4だけ全次数が下がる．
従って最高候補次数の項は
$(bc-ad)X^{a+c-1}\Xi^{b+d-1}$ となる．
これは右辺を単項式について計算したものに等しい．
一般の斉次多項式については線形性から，
非斉次の場合は最高次数成分だけが候補次数に
寄与することから結論が従う．
\end{proof}
\noindent\textbf{到達確認．}
$P=\partial^2+x$，$Q=x^3$ とする．
$PQ$ と $[P,Q]$ の正規形を求め，
交換子の全次数が上の評価と整合することを確かめよ．

\par\noindent\textbf{解答．}
交換公式から
\[
 \partial^2x^3=x^3\partial^2+6x^2\partial+6x
\]
であるから
\[
 \begin{aligned}
 PQ&=x^3\partial^2+6x^2\partial+6x+x^4,\\
 QP&=x^3\partial^2+x^4,\\
 [P,Q]&=6x^2\partial+6x.
 \end{aligned}
\]
$P$ の全次数は2，$Q$ は3であり，
交換子の全次数は3である．
これは $2+3-2=3$ に一致する．

\end{document}
